% Generated draft: requires mathematical and source review.
\documentclass[11pt]{article}
\usepackage[margin=25mm]{geometry}
\usepackage{fontspec}
\setmainfont{Latin Modern Roman}
\usepackage{amsmath,amssymb,mathtools,unicode-math}
\setmathfont{Latin Modern Math}
\usepackage{enumitem,needspace,xurl,hyperref}
\setlength{\parindent}{0pt}
\setlength{\parskip}{5pt}
\setlength{\emergencystretch}{3em}
\newcommand{\N}{\mathbb{N}}
\newcommand{\Z}{\mathbb{Z}}
\newcommand{\Q}{\mathbb{Q}}
\newcommand{\R}{\mathbb{R}}
\newcommand{\C}{\mathbb{C}}
\newcommand{\F}{\mathbb{F}}
\newcommand{\A}{\mathbb{A}}
\newcommand{\PP}{\mathbb{P}}
\newcommand{\E}{\mathbb{E}}
\newcommand{\eps}{\varepsilon}
\newcommand{\dd}{\,\mathrm{d}}
\begin{document}
\section*{UnsolvedMath ID: 10900025}
\textbf{Source problem code:} catalogue entry
\subsection{Definitions and mathematical statement}
For a finite simplicial graph \(\Lambda\), the associated right-angled Coxeter group is
\[W_\Lambda=\langle s_v\ (v\in V\Lambda)\mid s_v^2=1,\ [s_v,s_w]=1\text{ when }\{v,w\}\in E\Lambda\rangle.\]
The intended finite-generation convention is supported by the summary's number of generators but is not explicit in the short question. An affine action on \(\R^n\) has the form
\[g\cdot x=L(g)x+b(g),\]
where \(L:W_\Lambda\to\operatorname{GL}_n(\R)\) is a representation and
\[b(gh)=b(g)+L(g)b(h).\]
It is proper if for every compact \(C\subset\R^n\), the set \(\{g:gC\cap C\ne\varnothing\}\) is finite. Properness allows finite stabilizers and does not mean a free action; no cocompactness or preservation of a Euclidean metric is required.

Kassel and Mann ask to determine, for each given \(W_\Lambda\),
\[a(W_\Lambda)=\min\{n:W_\Lambda\text{ admits a proper affine action on }\R^n\},\]
with \(+\infty\) if there is no such dimension. Equivalently, find the sharp dimension admitting some pair \((L,b)\) with the above properness, and exclude every smaller allowed dimension. The convention whether \(n=0\) is allowed matters for finite groups and is not specified. The task asks for the minimum for a general given group, rather than merely a uniform upper bound or existence in some dimension.

The supplied summary attributes proper-action existence and dimension reductions to Danciger--Gu\'eritaud--Kassel. It specifically reports a \(k\)-dimensional bound, written \(\R^{p+q+1}\) with \(p+q+1=k\), for \(k\) generators, and improvements for particular cohomological-dimension-two and -four examples. The exact dimension formula is a supplied claim requiring verification from that paper; a standard linear representation's dimension alone does not establish an affine proper-action dimension. The summary expressly says that these results do not determine the sharp minimum for general right-angled Coxeter groups.
\subsection{Short English statement}
Determine the smallest affine dimension in which a given right-angled Coxeter group acts properly. Existing existence and upper-bound claims in the record are distinct from this sharp minimization problem.
\subsection{Sources}
\begin{itemize}
\item Local UnsolvedMath registry, permanent ID \(10900025\), exact source code \texttt{catalogue entry}.
\item Kassel and Mann, attribution in Delp, Hoffoss and Manning, \emph{Problems In Groups, Geometry, and Three-Manifolds} (2015), subsection 4.6, PDF page 3: \url{https://arxiv.org/pdf/1512.04620#page=3}; abstract \url{https://arxiv.org/abs/1512.04620}. Registry access 2026-07-29; source-tex extraction recorded as checked against rendered PDF text.
\item Danciger--Gu\'eritaud--Kassel, \emph{Proper affine actions for right-angled Coxeter groups}, \url{https://arxiv.org/abs/1804.03132}, reported by the summary as published in 2020. Publication details and the exact dimension statement are not independently checked here.
\end{itemize}
\end{document}
